% Editable overview. Add study photographs at the paths below.
% Keep each image's case ID, video timestamp and instrument group in its source notes.
% Author confirmed that the existing start annotations mark needle pickup.
\begin{figure*}[t]
\centering
\begingroup
% Target photo aspect ratio: width:height = 12:13 (960 x 1040 source frames).
\newlength{\overviewphotowidth}
\setlength{\overviewphotowidth}{.18\linewidth}
\newlength{\overviewimagewidth}
\newlength{\overviewimageheight}
% Arguments: card width, visible stage label, image path, photo-selection note.
% The fourth argument documents the intended image without adding a second caption.
\newcommand{\overviewstage}[4]{%
 \begin{minipage}[c]{#1}\centering
  \setlength{\overviewimagewidth}{\overviewphotowidth}%
  \ifdim\overviewimagewidth>\linewidth
   \setlength{\overviewimagewidth}{\linewidth}%
  \fi
  \setlength{\overviewimageheight}{\dimexpr\overviewimagewidth*13/12\relax}%
  \parbox[c][\overviewimageheight][c]{\overviewimagewidth}{\centering
   \IfFileExists{#3}{%
    \includegraphics[width=\overviewimagewidth,height=\overviewimageheight,keepaspectratio]{#3}%
   }{%
    \includegraphics[width=\overviewimagewidth,height=\overviewimageheight]{example-image-duck-portrait}%
   }%
  }\par\vspace{2pt}%
  \parbox[c][.28in][c]{\linewidth}{\centering\small\bfseries #2}%
 \end{minipage}}

\makebox[\linewidth][l]{\textbf{A. Recorded procedure}}\par\smallskip
\overviewstage{.18\linewidth}{Dissection}{figures/overview/dissection.jpg}{Dissection and tool exposure}%
\hspace{1pt}\raisebox{12pt}{\textcolor{gray}{$\rightarrow$}}\hspace{1pt}%
\overviewstage{.18\linewidth}{}{figures/overview/muscle_repair.jpg}{Tool use during muscle repair}%
\hspace{1pt}%
\overviewstage{.18\linewidth}{Repair phases}{figures/overview/soft_palate.jpg}{Tool use during soft-palate repair}%
\hspace{1pt}%
\overviewstage{.18\linewidth}{}{figures/overview/hard_palate.jpg}{Tool use during hard-palate repair}%
\hspace{1pt}\raisebox{12pt}{\textcolor{gray}{$\rightarrow$}}\hspace{1pt}%
\overviewstage{.18\linewidth}{Closure}{figures/overview/closure.jpg}{Completed closure}%
\par\smallskip

\makebox[\linewidth][l]{\textbf{B. Suturing annotation}}\par\smallskip
\overviewstage{.18\linewidth}{Needle pickup}{figures/overview/needle_pickup.jpg}{Needle grasp and orientation}%
\hspace{1pt}\raisebox{12pt}{\textcolor{gray}{$\rightarrow$}}\hspace{1pt}%
\overviewstage{.18\linewidth}{Needle exits tissue}{figures/overview/needle_exit.jpg}{Needle passage and tissue handling}%
\hspace{1pt}\raisebox{12pt}{\textcolor{gray}{$\rightarrow$}}\hspace{1pt}%
\overviewstage{.18\linewidth}{Last throw tightened}{figures/overview/knot_tightened.jpg}{Knot tightening}%
\hspace{1pt}\raisebox{12pt}{\textcolor{gray}{$\rightarrow$}}\hspace{1pt}%
\overviewstage{.18\linewidth}{Suture cutting and\\thread removal}{figures/overview/assistant_cut.jpg}{Assistant cutting and removal}%
% Timing endpoints: needle pickup, final throw tightening, and suture cutting.
\endgroup
\caption{Procedure overview. (A) Representative dissection, repair, and closure stages. (B) Suturing events. Suture placement time runs from needle pickup to final throw tightening; the subsequent interval ends at suture cutting. Thread removal is shown for context.}
\label{fig:workflow}
\end{figure*}
